\begin{table}
\caption{Tabela 5. Comparativo (D-027, emendado na D-029). Fração das células com p ≤ 0,05 contra as substitutas em 2001–2024, por variável, ao lado do limite superior do IC 95% da taxa de rejeição na nula realista mais permissiva ajustada à mesma variável (capítulo 5). n′: número efetivo de células independentes. ▲: acima do limite mesmo com o IC da fração; △: acima só pontualmente; ≈: a nula rejeita em mais de 95% e a comparação não informa. *: na subamostra de D-024. S-map: teste de Diebold–Mariano do ganho sobre a previsão linear, horizonte de 1 dia. Limite e leitura por célula: a mesma regra com as quatro nulas ajustadas a cada célula da subamostra (39 substitutas por série), Xingu / Tapajós; ·: dentro do limite; –: estimador sem nula por célula.}
\label{tab:07_descritores::natureza_nulas_realistas}
\begin{tabular}{llllllll}
\toprule
Variável & Descritor & Nula mais permissiva & Limite & Xingu & Tapajós & Limite por célula & Leitura por célula \\
\midrule
temperatura & Entropia de permutação & AR(1), memória sazonal & 35,6\% & 74\% △ (n′ 4) & 52\% △ (n′ 3) & 56,8\% & △ / · \\
temperatura & Complexidade estatística & AR(1), memória sazonal & 29,8\% & 69\% △ (n′ 4) & 51\% △ (n′ 3) & 53,5\% & △ / · \\
temperatura & Erro de previsão não linear & AR(1)-GARCH & 32,8\% & 0\% (n′ 6) & 0\% (n′ 24) & 23,2\% & · / · \\
temperatura & RQA · determinismo & AR(1)-GARCH & 100,0\% & 100\% ≈ (n′ 8) & 100\% ≈ (n′ 16) & 100,0\% & ≈ / ≈ \\
temperatura & RQA · laminaridade & AR(1)-GARCH & 100,0\% & 100\% ≈ (n′ 9) & 100\% ≈ (n′ 13) & 88,3\% & △ / △ \\
temperatura & Dimensão de correlação & AR(1)-GARCH & 63,9\% & 100\% ▲ (n′ 12) & 100\% ▲ (n′ 22) & 48,6\% & ▲ / ▲ \\
temperatura & Assimetria temporal & SETAR, dois regimes & 46,5\% & 100\% ▲ (n′ 11) & 100\% ▲ (n′ 13) & 46,9\% & ▲ / ▲ \\
temperatura & Lyapunov (Rosenstein) & AR(1)-GARCH & 87,3\% & 14\% (n′ 7) & 35\% (n′ 9) & 41,0\% & · / · \\
temperatura & Lyapunov (Wolf)* & AR(1)-GARCH & 41,2\% & 43\% △ (n′ 15) & 37\% (n′ 18) & – & – \\
temperatura & Lyapunov (Kantz)* & AR(1)-GARCH & 99,4\% & 98\% ≈ (n′ 6) & 100\% ≈ (n′ 8) & – & – \\
temperatura & Cao · D\_E2* & AR(1) com tendência & 20,4\% & 2\% (n′ 60) & 18\% (n′ 48) & – & – \\
temperatura & S-map (ganho, 1 dia) & SETAR, dois regimes & 56,7\% & 69\% △ (n′ 8) & 86\% △ (n′ 13) & 43,6\% & △ / ▲ \\
déficit de pressão de vapor & Entropia de permutação & AR(1), memória sazonal & 99,9\% & 81\% ≈ (n′ 4) & 93\% ≈ (n′ 6) & 94,8\% & · / · \\
déficit de pressão de vapor & Complexidade estatística & AR(1), memória sazonal & 98,4\% & 77\% ≈ (n′ 4) & 90\% ≈ (n′ 6) & 94,2\% & · / · \\
déficit de pressão de vapor & Erro de previsão não linear & AR(1), memória sazonal & 35,6\% & 0\% (n′ 7) & 0\% (n′ 38) & 45,3\% & · / · \\
déficit de pressão de vapor & RQA · determinismo & AR(1), memória sazonal & 100,0\% & 100\% ≈ (n′ 6) & 100\% ≈ (n′ 7) & 100,0\% & ≈ / ≈ \\
déficit de pressão de vapor & RQA · laminaridade & AR(1), memória sazonal & 100,0\% & 100\% ≈ (n′ 5) & 100\% ≈ (n′ 7) & 99,8\% & ≈ / ≈ \\
déficit de pressão de vapor & Dimensão de correlação & AR(1), memória sazonal & 99,9\% & 100\% ≈ (n′ 9) & 99\% ≈ (n′ 11) & 89,7\% & △ / △ \\
déficit de pressão de vapor & Assimetria temporal & SETAR, dois regimes & 61,5\% & 72\% △ (n′ 5) & 90\% △ (n′ 8) & 75,4\% & · / △ \\
déficit de pressão de vapor & Lyapunov (Rosenstein) & AR(1)-GARCH & 95,8\% & 33\% ≈ (n′ 9) & 48\% ≈ (n′ 33) & 89,0\% & · / · \\
déficit de pressão de vapor & Lyapunov (Wolf)* & AR(1), memória sazonal & 79,4\% & 50\% (n′ 14) & 53\% (n′ 8) & – & – \\
déficit de pressão de vapor & Lyapunov (Kantz)* & AR(1)-GARCH & 100,0\% & 72\% ≈ (n′ 4) & 56\% ≈ (n′ 17) & – & – \\
déficit de pressão de vapor & Cao · D\_E2* & AR(1) com tendência & 20,4\% & 12\% (n′ 60) & 19\% (n′ 15) & – & – \\
déficit de pressão de vapor & S-map (ganho, 1 dia) & SETAR, dois regimes & 73,0\% & 42\% (n′ 12) & 69\% (n′ 8) & 67,1\% & · / △ \\
\bottomrule
\end{tabular}
\end{table}
